\documentclass[UTF8,a4paper,11pt]{ctexart}

\usepackage[margin=2.2cm]{geometry}
\usepackage{booktabs}
\usepackage{array}
\usepackage{longtable}
\usepackage{hyperref}
\usepackage{xcolor}
\usepackage{enumitem}
\usepackage{graphicx}
\usepackage{float}
\usepackage{tikz}
\usetikzlibrary{arrows.meta, positioning, shapes.geometric, fit, backgrounds}

\tikzset{
  archbox/.style={
    rectangle,
    rounded corners=2pt,
    draw=black,
    fill=gray!8,
    align=center,
    minimum width=3.5cm,
    minimum height=0.9cm,
    font=\small
  },
  storebox/.style={
    cylinder,
    shape border rotate=90,
    aspect=0.25,
    draw=black,
    fill=blue!8,
    align=center,
    minimum height=1.15cm,
    minimum width=3.6cm,
    font=\small
  },
  phasebox/.style={
    rounded corners=4pt,
    draw=gray!65,
    dashed,
    inner sep=0.35cm
  },
  archarrow/.style={-{Stealth[length=2.4mm]}, thick},
  tracearrow/.style={-{Stealth[length=2.2mm]}, dashed, gray!70}
}

\hypersetup{
  colorlinks=true,
  linkcolor=black,
  urlcolor=blue,
  citecolor=black
}

\setlist[itemize]{leftmargin=2em,itemsep=0.2em}
\setlist[enumerate]{leftmargin=2em,itemsep=0.2em}
\setlength{\emergencystretch}{3em}
\setlength{\tabcolsep}{4pt}
\renewcommand{\arraystretch}{1.18}
\newcolumntype{L}[1]{>{\raggedright\arraybackslash}p{#1}}

\title{\textbf{长期记忆对话 Agent 中期进展报告}}
\author{汤岑宸 \quad 学号：231880523 \quad 曹信阳 \quad 学号：231880114}
\date{2026年6月}

\begin{document}
\maketitle

\section{项目目标}

本项目旨在实现一个支持长期记忆的对话 Agent，使其能够在跨会话、跨时间的交互中保存并利用历史信息。系统需要从多轮、多会话对话中抽取结构化记忆，在回答新问题时检索相关记忆，并通过记忆更新、冲突处理或遗忘机制提升长期对话问答能力。

与普通上下文增强方法不同，本项目重点不只是把原始对话直接拼入 prompt，而是将原始对话日志加工为可管理、可追踪、可检索的记忆单元。最终系统将在基于 LoCoMo 的长期对话评测集上进行评估，并与 No-memory、Full-context 和 Vanilla RAG 等基线方法进行对比。

\section{当前进度}

截至中期提交阶段，项目已完成以下工作：

\begin{enumerate}
  \item 完成任务需求分析，明确长期记忆 Agent 的核心目标，包括记忆写入、记忆存储、记忆检索、记忆更新和回答生成。
  \item 阅读并整理课程提供的项目说明与评测工具包，明确统一实验流程：准备评测集、运行 Agent 生成预测、使用 LLM-as-Judge 评测结果。
  \item 完成 Python 虚拟环境配置，安装 PyTorch、sentence-transformers、OpenAI SDK 和 vLLM。当前机器可识别 NVIDIA GeForce RTX 4060 Laptop GPU，并可运行 CUDA 推理。
  \item 使用本地下载的 \texttt{Qwen2.5-3B-Instruct-AWQ} 通过 vLLM 提供 OpenAI 兼容接口作为生成模型。Vanilla RAG 使用本地 \texttt{BAAI/bge-m3} embedding 模型进行向量检索，embedding 计算设置为 CPU，LLM 推理使用 GPU。
  \item 使用助教脚本 \texttt{prepare\_eval\_set.py} 生成 200 问题评测集，并完成 No-memory、Full-context、Vanilla RAG 三组 baseline 以及当前自定义 MemoryAgent 的预测生成。
  \item 使用 DeepSeek 线上 API 和助教脚本 \texttt{run\_judge.py} 完成 LLM-as-Judge 正式测评，获得当前阶段的主指标、调用成本和端到端响应时间。
\end{enumerate}

需要说明的是，当前工具包本地缓存数据源为 \texttt{locomo10.json}，因此本阶段生成的评测集包含 10 段可用对话、200 个问题，并保证四类问题各 50 个。

\section{系统方案设计}

当前实现的长期记忆 Agent 采用“写入--更新--存储--检索--生成”的流水线结构。整体目标不是直接把全部原始对话塞入 prompt，而是在 \texttt{ingest} 阶段先把历史对话转化为可管理的长期记忆单元，再在 \texttt{answer} 阶段检索相关记忆辅助回答。

\begin{figure}[H]
\centering
\resizebox{0.98\linewidth}{!}{
\begin{tikzpicture}[x=1cm,y=1cm]

\node[font=\small\bfseries] at (0,1.0) {写入阶段};
\node[font=\small\bfseries] at (5.2,1.0) {记忆库};
\node[font=\small\bfseries] at (10.4,1.0) {回答阶段};

\node[archbox, minimum width=3.2cm] (conv) at (0,0) {Conversation\\sessions};
\node[archbox, minimum width=3.2cm] (writer) at (0,-1.6) {Memory Writer\\LLM 事实抽取};
\node[archbox, minimum width=3.2cm] (updater) at (0,-3.2) {Memory Updater\\去重 / 更新};

\node[storebox, minimum width=3.5cm] (store) at (5.2,-1.6) {Memory Store\\metadata + bge-m3 index};

\node[archbox, minimum width=3.2cm] (query) at (10.4,0) {Question};
\node[archbox, minimum width=3.2cm] (retriever) at (10.4,-1.6) {Memory Retriever\\三因子排序};
\node[archbox, minimum width=3.2cm] (controller) at (10.4,-3.2) {Agent Controller\\prompt 组织};
\node[archbox, minimum width=3.2cm] (answer) at (10.4,-4.8) {LLM Answer};

\node[archbox, minimum width=5.8cm, fill=yellow!10] (trace) at (5.2,-6) {Trace / Logs：检索记忆、三因子分数、prompt、输出};

\draw[archarrow] (conv) -- node[left, font=\scriptsize]{ingest} (writer);
\draw[archarrow] (writer) -- (updater);
\draw[archarrow] (updater.east) -- node[above, font=\scriptsize]{clean memories} (store.west);

\draw[archarrow] (query) -- node[right, font=\scriptsize]{answer} (retriever);
\draw[archarrow] (store.east) -- node[above, font=\scriptsize]{top-k} (retriever.west);
\draw[archarrow] (retriever) -- (controller);
\draw[archarrow] (controller) -- (answer);

\draw[tracearrow] (writer.south east) .. controls (2.1,-3.9) and (3.2,-4.7) .. (trace.north west);
\draw[tracearrow] (retriever.south west) .. controls (8.1,-3.7) and (7.1,-4.6) .. (trace.north east);
\draw[tracearrow] (controller.south west) .. controls (8.8,-4.7) and (7.5,-5.1) .. (trace.east);

\end{tikzpicture}
}
\caption{自定义 MemoryAgent 的整体架构}
\label{fig:memory-agent-architecture}
\end{figure}

其中，Memory Writer 和 Memory Updater 主要参考 Mem0 的“记忆提取--冲突检测--更新”流程；Memory Retriever 参考 Generative Agents 的 relevance、recency 和 importance 三因子检索；Memory Store 中的时间衰减是对 MemoryBank 遗忘思想的轻量近似。当前版本优先保证结构完整、日志可追踪和实验可消融，而不是追求复杂度。

\begin{longtable}{L{2.6cm}L{4.0cm}L{6.0cm}}
\caption{自定义 MemoryAgent 架构}\\
\toprule
模块 & 输入 / 输出 & 当前实现 \\
\midrule
\endfirsthead
\toprule
模块 & 输入 / 输出 & 当前实现 \\
\midrule
\endhead
\midrule
\endfoot
\bottomrule
\endlastfoot
Memory Writer & 输入原始 session，输出结构化记忆 & 使用本地 Qwen 对每个 session 分块抽取 JSON 记忆，字段包括 text、subject、attribute、importance、source、timestamp \\
Memory Updater & 输入候选记忆，输出清洗后的记忆 & 文本级去重；对相同 subject + attribute 的记忆保留较新的版本 \\
Memory Store & 输入记忆单元，输出可检索索引 & 使用列表保存元数据，用 \texttt{BAAI/bge-m3} 编码记忆文本，矩阵乘法计算相似度 \\
Memory Retriever & 输入 query，输出 top-k 记忆 & 使用 relevance、recency、importance 三因子加权排序，默认权重为 0.70、0.15、0.15 \\
Agent Controller & 输入问题，输出最终答案 & 串联写入、更新、存储、检索和生成；记录检索结果、prompt、三因子分数和 LLM 输出 \\
\end{longtable}

\subsection{Memory Writer}

Memory Writer 负责从原始多会话对话中抽取结构化记忆。当前版本参考 Mem0 的工程流程，使用本地 \texttt{Qwen2.5-3B-Instruct-AWQ} 对每个 session 分块抽取 JSON 形式的记忆单元。每条记忆包含：

\begin{itemize}
  \item \texttt{text}：一条独立、简短的事实性记忆；
  \item \texttt{subject}：记忆涉及的人或实体；
  \item \texttt{attribute}：记忆类型，如 preference、location、job、event、relationship、plan、update 等；
  \item \texttt{importance}：1--5 的重要性分数；
  \item \texttt{source} 和 \texttt{timestamp}：用于追踪来源和时间。
\end{itemize}

该模块避免直接把原始对话切片作为最终记忆，而是尝试将对话压缩为更清晰的事实性记忆。这样做的优点是记忆更短、更容易管理；缺点是如果 Writer 漏掉了关键事实，后续检索和生成都无法恢复该信息。

\subsection{Memory Store}

Memory Store 负责保存结构化记忆，并为检索提供索引。当前版本使用列表结构保存元数据，同时使用本地 \texttt{BAAI/bge-m3} 生成向量表示，并通过矩阵乘法计算 query 与记忆之间的余弦相似度。每条记忆还保存插入顺序、访问次数、时间戳和重要性等字段，方便后续加入遗忘曲线、访问增强或 FAISS / Chroma 索引。

\subsection{Memory Retriever}

Memory Retriever 在回答问题时根据 query 检索相关记忆。当前版本参考 Generative Agents，引入 relevance、recency 和 importance 三因子排序：

\[
\mathrm{Score}(m, q)
= 0.70 \cdot \mathrm{Relevance}(m, q)
+ 0.15 \cdot \mathrm{Recency}(m)
+ 0.15 \cdot \mathrm{Importance}(m)
\]

其中，Relevance 来自 \texttt{bge-m3} embedding 相似度；Recency 使用时间衰减近似，越新的记忆分数越高；Importance 来自 Writer 抽取阶段给出的 1--5 分并归一化。当前默认返回 top-6 条记忆，并在 trace 中保存每条记忆的总分和三个子分数，方便后续分析检索错误。

\subsection{Memory Updater}

记忆管理机制当前实现了轻量去重和更新。当两条记忆文本规范化后相同时，只保留一条并保留较高 importance；当两条记忆具有相同的 \texttt{subject + attribute} 时，系统保留时间更新的记忆。该机制对应 Mem0 中的冲突检测与更新思想，但当前没有额外调用 LLM 判别冲突，因此成本较低。

\subsection{Agent Controller}

Agent Controller 负责组织完整流程。对于每段对话，评测脚本先调用 \texttt{ingest(conversation)}，此时系统完成记忆写入、更新和索引构建；随后对该对话下的每个问题调用 \texttt{answer(question)}，Controller 检索相关记忆并构造回答 prompt。

当前回答 prompt 只包含检索出的结构化记忆，而不直接包含原始对话片段。这样可以测试“结构化长期记忆”本身是否有效，但也带来了当前 Ours 效果较差的核心风险：一旦 Writer 抽取时丢失细节，Generator 无法从原始对话中补救。因此，下一步计划加入“结构化记忆 + 原始证据片段”的回退机制。

\section{初步实验设置}

本阶段实验使用助教提供的统一脚本和接口规范。生成阶段使用本地 vLLM 部署的 \texttt{Qwen2.5-3B-Instruct-AWQ}，Judge 阶段使用 DeepSeek 线上 API 的 \texttt{deepseek-v4-flash}。Vanilla RAG 的检索模型为本地 \texttt{BAAI/bge-m3}。所有 baseline 使用同一底座生成模型，变化只发生在记忆使用方式上。

\begin{longtable}{L{2.8cm}L{4.2cm}L{6.0cm}}
\caption{当前方法设置}\\
\toprule
方法 & 记忆使用方式 & 说明 \\
\midrule
\endfirsthead
\toprule
方法 & 记忆使用方式 & 说明 \\
\midrule
\endhead
\midrule
\endfoot
\bottomrule
\endlastfoot
No-memory & 不使用历史对话 & 仅将当前问题输入 LLM \\
Full-context & 截断后的历史上下文 & 将对话历史拼入 prompt \\
Vanilla RAG & 原始对话切片检索 & 使用向量检索取 top-k 片段辅助回答 \\
Ours & 结构化记忆检索 & LLM 写入记忆，三因子检索，轻量更新 \\
\end{longtable}

评测指标包括三类：第一，QA 准确率，即助教脚本基于 LLM-as-Judge 计算的 \texttt{Judge Score}，这是本项目主指标；第二，平均每题 LLM 调用次数，用于衡量系统调用成本；第三，平均端到端响应时间，用于衡量系统效率。

\section{正式测评结果}

表 \ref{tab:judge-main} 给出了当前 200 问题评测集上的正式 LLM-as-Judge 结果。Judge 输出标签包括 \texttt{CORRECT}、\texttt{PARTIAL} 和 \texttt{WRONG}，对应分数分别为 1.0、0.5 和 0.0，总体得分为所有题目得分的平均值。

\begin{longtable}{lccccc}
\caption{当前方法的正式 Judge 结果}
\label{tab:judge-main}\\
\toprule
方法 & Judge Score & F1 & EM & 平均 LLM 调用 & 平均端到端时间 \\
\midrule
\endfirsthead
\toprule
方法 & Judge Score & F1 & EM & 平均 LLM 调用 & 平均端到端时间 \\
\midrule
\endhead
\midrule
\endfoot
\bottomrule
\endlastfoot
No-memory & 0.0000 & 0.0090 & 0.000 & 1.0 & 0.064s \\
Full-context & 0.0500 & 0.0547 & 0.005 & 1.0 & 0.865s \\
Vanilla RAG & \textbf{0.1300} & \textbf{0.1455} & \textbf{0.015} & 1.0 & 3.269s \\
Ours & 0.0275 & 0.0653 & 0.010 & 2.69 & 12.785s \\
\end{longtable}

从总体结果看，Vanilla RAG 的 Judge Score 最高，为 0.1300；Full-context 为 0.0500；No-memory 为 0.0000；当前自定义 MemoryAgent 仅为 0.0275。该结果说明，长期记忆任务中历史信息非常关键，而直接检索原始对话片段的 Vanilla RAG 在当前阶段反而最稳。我们的系统虽然结构更完整，但因为 Memory Writer 会先压缩原始对话，若抽取阶段丢失关键信息，后续三因子检索无法恢复这些证据。同时，Ours 的平均每题 LLM 调用次数达到 2.69，端到端时间为 12.785s，说明当前结构化记忆写入带来了明显成本，但尚未带来准确率收益。

\begin{longtable}{lcccc}
\caption{按问题类型拆分的 Judge Score}
\label{tab:judge-category}\\
\toprule
方法 & multi-hop & open-domain & single-hop & temporal \\
\midrule
\endfirsthead
\toprule
方法 & multi-hop & open-domain & single-hop & temporal \\
\midrule
\endhead
\midrule
\endfoot
\bottomrule
\endlastfoot
No-memory & 0.000 & 0.000 & 0.000 & 0.000 \\
Full-context & 0.000 & 0.180 & 0.020 & 0.000 \\
Vanilla RAG & \textbf{0.040} & \textbf{0.270} & \textbf{0.180} & \textbf{0.030} \\
Ours & 0.000 & 0.060 & 0.050 & 0.000 \\
\end{longtable}

分类结果显示，Vanilla RAG 在 open-domain 和 single-hop 问题上提升最明显，分别达到 0.270 和 0.180；但所有方法在 multi-hop 和 temporal 上仍然较弱。这与长期记忆任务的难点一致：时间推理需要显式建模事件发生时间、更新时间和先后关系；多跳问题则要求系统跨多个会话片段整合信息。Ours 在 temporal 和 multi-hop 上均为 0，说明当前结构化记忆抽取没有稳定保留多跳证据和时间线信息。

\section{当前 MemoryAgent 效果较差的原因分析}

当前自定义 MemoryAgent 的效果明显低于 Vanilla RAG，主要原因并不是检索模型 \texttt{bge-m3} 本身，而是结构化记忆流水线中存在信息瓶颈。

\begin{enumerate}
  \item \textbf{写入阶段信息丢失。} Vanilla RAG 直接保留原始对话片段，而 Ours 先让本地 3B 模型把 session 压缩为少量记忆句子。长期记忆问答中的答案往往依赖细节，例如具体日期、事件对象、计划变化或一句短回复；这些细节在抽取时容易被概括掉。一旦 Writer 没有写入某个事实，Retriever 和 Generator 后面都无法恢复。
  \item \textbf{本地小模型不擅长稳定 JSON 记忆抽取。} \texttt{Qwen2.5-3B-Instruct-AWQ} 适合作为轻量生成模型，但用于长 session 的结构化事实抽取时，容易出现过度总结、遗漏实体、attribute 标注粗糙等问题。这会导致 Store 中的记忆看起来合理，但不一定覆盖评测问题所需答案。
  \item \textbf{三因子检索的 recency 和 importance 目前偏粗。} 当前所有高层记忆往往 importance 都较高，importance 区分度不足；时间戳解析和 recency 衰减也只是近似实现。这样三因子排序并没有充分发挥 Generative Agents 中的优势，反而可能把较新的泛化记忆排在真正相关的旧事实前面。
  \item \textbf{成本高但有效证据少。} Ours 在 ingest 阶段需要额外调用 LLM 抽取记忆，因此平均每题 LLM 调用次数达到 2.69，端到端时间也显著增加。但这些额外调用主要用于压缩历史，而不是直接回答问题；如果压缩质量不足，就会形成“更贵但更少证据”的系统。
\end{enumerate}

因此，当前结果说明：系统架构已经跑通，但 Memory Writer 还不是一个可靠的信息保留模块。后续改进应优先降低写入损失，而不是继续增加复杂检索权重。

\section{结果水平与阶段性判断}

结合当前任务和 LoCoMo 风格长期记忆问答的难度，可以将 Judge Score 粗略理解为以下几个区间：

\begin{longtable}{L{3.0cm}L{9.5cm}}
\caption{Judge Score 的经验性解释}\\
\toprule
Judge Score 区间 & 效果判断 \\
\midrule
\endfirsthead
\toprule
Judge Score 区间 & 效果判断 \\
\midrule
\endhead
\midrule
\endfoot
\bottomrule
\endlastfoot
0.00--0.05 & 基本未解决长期记忆问题，接近无记忆或随机回答 \\
0.05--0.12 & baseline 水平，能够答中少量简单题 \\
0.12--0.25 & 有一定长期记忆能力，检索或记忆模块开始有效 \\
0.25--0.40 & 效果较好，能够作为明显改进结果 \\
0.40 以上 & 效果较强，通常需要高质量记忆写入、检索排序和更新机制 \\
\end{longtable}

据此判断，当前 Vanilla RAG 的 0.1300 已经达到“有一定长期记忆能力”的初步区间，而 Ours 的 0.0275 仍接近无效系统。中期阶段该结果的意义主要在于：实验环境、数据准备、生成脚本、线上 Judge 流程和自定义 MemoryAgent 架构都已经完整跑通；同时结果也暴露出，结构化记忆并不会天然优于原始 RAG，关键在于 Writer 是否能高保真保留答案证据。后续自定义 MemoryAgent 应首先追平并超过 Vanilla RAG（0.1300），再进一步做三因子检索和更新机制的消融。

\section{下一步计划}

后续将围绕当前 MemoryAgent 的失败原因进行改进，重点包括：

\begin{enumerate}
  \item 改进 Memory Writer，不再只输出高度压缩的摘要记忆，而是保留更多原子事实，尤其是人物、时间、地点、计划和偏好变化。
  \item 设计“原始证据回退”机制：检索结构化记忆的同时保留对应原始对话片段，回答时同时提供 memory 和 source snippets，减少写入阶段造成的信息丢失。
  \item 做检索策略消融，对比只使用 relevance、relevance + recency、relevance + recency + importance，判断三因子权重是否真的有效。
  \item 改进时间处理，对 temporal 问题显式保留事件时间和更新顺序，而不是仅依赖 embedding 相似度。
  \item 基于 trace 挑选 3--5 个 bad case，分别判断错误来自写入丢失、检索未命中还是生成未利用。
\end{enumerate}

\section{风险与解决方案}

当前主要风险包括三类。第一，长期对话数据较长，Full-context 容易触发上下文长度限制或显存压力。解决方案是保留 vLLM 8K 上下文配置，并在 baseline 中使用固定截断策略。第二，LLM-as-Judge 存在 API 成本和速率限制。解决方案是先用小样本调试，正式评测时使用 DeepSeek API 并控制并发数。第三，结构化记忆抽取质量可能不稳定。解决方案是在小样本上人工检查记忆单元，并通过 bad case 将错误归因到写入失败、检索失败或生成失败。

\section{总结}

中期阶段已完成长期记忆 Agent 的需求分析、系统方案设计、环境配置、baseline 生成、自定义 MemoryAgent 实现和线上 LLM-as-Judge 测评。实验结果表明，Vanilla RAG 当前明显优于 No-memory、Full-context 和 Ours，说明在当前数据和本地 3B 生成模型设置下，直接保留原始对话片段比低质量结构化抽取更可靠。我们的 MemoryAgent 效果较差，主要原因是写入阶段压缩导致关键信息丢失，同时带来了额外 LLM 调用和延迟。下一阶段将重点改进 Memory Writer 的高保真事实保留能力，并加入原始证据回退机制，再通过消融实验验证三因子检索和更新机制是否真正带来增益。

\end{document}
